% GENERATED FILE. Edit canonical lessons, then run npm run generate:books.
% canonical-source-hash: fnv1a64:9247873e009e873d
% canonical-lessons: RU-C497-vospitaniye, RU-C497-balovat, RU-C497-slushatsya, RU-C497-shalit, RU-C497-neposlushnyy

\chapter{Bringing up Children}
\label{ch:ru-bringing-up-children}
\hlchaptermodality{497}

\begin{chapteropening}
\textbf{By the end of this chapter:} \emph{I can talk about bringing up children and how they behave.}

Complete the last lesson of chapter 497: I can talk about bringing up children and how they behave.
\end{chapteropening}

\section[\texorpdfstring{\ru{воспитание} (vospitániye)}{воспитание (vospitániye)}]{\ru{воспитание} (vospitániye) — upbringing}

\label{lesson:RU-C497-vospitaniye}



\begin{quote}
Before the new one: say the Russian for to lull to sleep, then the Russian for to swaddle.
\end{quote}

\subsection*{You'll want to know: \ru{воспитание}}
\textbf{\ru{воспитание}} (\emph{vospitániye}) — \textquotedblleft{}upbringing\textquotedblright{}.

\begin{quote}
\ru{Хорошее} \ru{воспитание} \ru{видно} \ru{сразу}.
\end{quote}

\begin{quote}
\emph{Khorósheye vospitániye vídno srázu.}
\end{quote}

\begin{quote}
A good upbringing shows at once.
\end{quote}

Neuter. The verbs are \emph{\ru{воспитывать} / \ru{воспитать}} — to bring up.

\begin{cousinweb}
\emph{\ru{воспитание}} is built on \emph{\ru{воспитать}} 'to bring up', from \emph{\ru{вос}-} and \emph{\ru{питать}} 'to feed, to nourish'; EtymDB has no entry for the noun.
\end{cousinweb}

\begin{grammarlens}[title={What you've built}]
One more word for telling, explaining and describing at B1.
\end{grammarlens}

\subsection*{Your turn}
\begin{itemize}
\raggedright
  \item \emph{Say it:} \emph{vospitániye}
  \item \emph{Say it:} \emph{vospitániye}, once more
  \item \emph{Say it:} a sentence of your own with \emph{vospitániye}
  \item \emph{Recall:} say the Russian for to lull to sleep, then the Russian for to swaddle, then say \emph{vospitániye} again
\end{itemize}

\begin{tcolorbox}[breakable,colback=teal!4,colframe=teal!35!black,title={Before you move on}]
What does \ru{воспитание} mean? (\textquotedblleft{}Upbringing\textquotedblright{}.) Say it once more.
\end{tcolorbox}
\section[\texorpdfstring{\ru{баловать} (balovát)}{баловать (balovát)}]{\ru{баловать} (balovát) — to spoil (a child)}

\label{lesson:RU-C497-balovat}



\begin{quote}
Before the new one: say the Russian for to swaddle, then the Russian for upbringing.
\end{quote}

\subsection*{You'll want to know: \ru{баловать}}
\textbf{\ru{баловать}} (\emph{balovát}) — \textquotedblleft{}to spoil (a child)\textquotedblright{}.

\begin{quote}
\ru{Бабушка} \ru{балует} \ru{внуков} \ru{сладостями}.
\end{quote}

\begin{quote}
\emph{Bábushka balúyet vnúkov sládostyami.}
\end{quote}

\begin{quote}
Grandma spoils her grandchildren with sweets.
\end{quote}

Imperfective; present \ru{балую}, \ru{балуешь} (-\ru{ова}- becomes -\ru{у}-). Perfective \ru{избаловать}.

\begin{cousinweb}
\emph{\ru{баловать}} is inherited from Old East Slavic \emph{balovati} 'to flatter, to heal', which EtymDB traces to Proto-Slavic \emph{badlovati}.
\end{cousinweb}

\begin{grammarlens}[title={What you've built}]
One more word for telling, explaining and describing at B1.
\end{grammarlens}

\subsection*{Your turn}
\begin{itemize}
\raggedright
  \item \emph{Say it:} \emph{balovát}
  \item \emph{Say it:} \emph{balovát}, once more
  \item \emph{Say it:} a sentence of your own with \emph{balovát}
  \item \emph{Recall:} say the Russian for to swaddle, then the Russian for upbringing, then say \emph{balovát} again
\end{itemize}

\begin{tcolorbox}[breakable,colback=teal!4,colframe=teal!35!black,title={Before you move on}]
What does \ru{баловать} mean? (\textquotedblleft{}To spoil (a child)\textquotedblright{}.) Say it once more.
\end{tcolorbox}
\section[\texorpdfstring{\ru{слушаться} (slúshatsya)}{слушаться (slúshatsya)}]{\ru{слушаться} (slúshatsya) — to obey}

\label{lesson:RU-C497-slushatsya}



\begin{quote}
Before the new one: say the Russian for upbringing, then the Russian for to spoil.
\end{quote}

\subsection*{You'll want to know: \ru{слушаться}}
\textbf{\ru{слушаться}} (\emph{slúshatsya}) — \textquotedblleft{}to obey\textquotedblright{}.

\begin{quote}
\ru{Дети} \ru{не} \ru{всегда} \ru{слушаются} \ru{родителей}.
\end{quote}

\begin{quote}
\emph{Déti ne vsegdá slúshayutsya rodíteley.}
\end{quote}

\begin{quote}
Children don't always do what their parents say.
\end{quote}

Imperfective; perfective \ru{послушаться}. Takes the genitive: \emph{\ru{слушаться} \ru{родителей}}.

\begin{cousinweb}
\emph{\ru{слушаться}} is the reflexive of \emph{\ru{слушать}} 'to listen': to obey is to listen to someone. EtymDB has no entry for the reflexive.
\end{cousinweb}

\begin{grammarlens}[title={What you've built}]
One more word for telling, explaining and describing at B1.
\end{grammarlens}

\subsection*{Your turn}
\begin{itemize}
\raggedright
  \item \emph{Say it:} \emph{slúshatsya}
  \item \emph{Say it:} \emph{slúshatsya}, once more
  \item \emph{Say it:} a sentence of your own with \emph{slúshatsya}
  \item \emph{Recall:} say the Russian for upbringing, then the Russian for to spoil, then say \emph{slúshatsya} again
\end{itemize}

\begin{tcolorbox}[breakable,colback=teal!4,colframe=teal!35!black,title={Before you move on}]
What does \ru{слушаться} mean? (\textquotedblleft{}To obey\textquotedblright{}.) Say it once more.
\end{tcolorbox}
\section[\texorpdfstring{\ru{шалить} (shalít)}{шалить (shalít)}]{\ru{шалить} (shalít) — to be naughty}

\label{lesson:RU-C497-shalit}



\begin{quote}
Before the new one: say the Russian for to spoil, then the Russian for to obey.
\end{quote}

\subsection*{You'll want to know: \ru{шалить}}
\textbf{\ru{шалить}} (\emph{shalít}) — \textquotedblleft{}to be naughty\textquotedblright{}.

\begin{quote}
\ru{На} \ru{перемене} \ru{мальчики} \ru{шалят} \ru{в} \ru{коридоре}.
\end{quote}

\begin{quote}
\emph{Na pereméne málchiki shalyát v koridóre.}
\end{quote}

\begin{quote}
At break the boys get up to mischief in the corridor.
\end{quote}

Imperfective; perfective \ru{пошалить}. The noun is \emph{\ru{шалость}} — a prank.

\begin{cousinweb}
EtymDB lists \emph{\ru{шалить}} without a source, and the origin is not settled.
\end{cousinweb}

\begin{grammarlens}[title={What you've built}]
One more word for telling, explaining and describing at B1.
\end{grammarlens}

\subsection*{Your turn}
\begin{itemize}
\raggedright
  \item \emph{Say it:} \emph{shalít}
  \item \emph{Say it:} \emph{shalít}, once more
  \item \emph{Say it:} a sentence of your own with \emph{shalít}
  \item \emph{Recall:} say the Russian for to spoil, then the Russian for to obey, then say \emph{shalít} again
\end{itemize}

\begin{tcolorbox}[breakable,colback=teal!4,colframe=teal!35!black,title={Before you move on}]
What does \ru{шалить} mean? (\textquotedblleft{}To be naughty\textquotedblright{}.) Say it once more.
\end{tcolorbox}
\section[\texorpdfstring{\ru{непослушный} (neposlúshnyy)}{непослушный (neposlúshnyy)}]{\ru{непослушный} (neposlúshnyy) — disobedient, naughty}

\label{lesson:RU-C497-neposlushnyy}



\begin{quote}
Before the new one: say the Russian for to obey, then the Russian for to be naughty.
\end{quote}

\subsection*{You'll want to know: \ru{непослушный}}
\textbf{\ru{непослушный}} (\emph{neposlúshnyy}) — \textquotedblleft{}disobedient, naughty\textquotedblright{}.

\begin{quote}
\ru{Непослушный} \ru{ребёнок} \ru{убежал} \ru{от} \ru{мамы}.
\end{quote}

\begin{quote}
\emph{Neposlúshnyy rebyónok ubezhál ot mámy.}
\end{quote}

\begin{quote}
The naughty child ran away from his mum.
\end{quote}

The opposite is \emph{\ru{послушный}} — obedient.

\begin{cousinweb}
\emph{\ru{непослушный}} is built on \emph{\ru{не}-} and \emph{\ru{послушный}} 'obedient', from \emph{\ru{слушать}} 'to listen'; EtymDB has no entry.
\end{cousinweb}

\begin{grammarlens}[title={What you've built}]
One more word for telling, explaining and describing at B1.
\end{grammarlens}

\subsection*{Your turn}
\begin{itemize}
\raggedright
  \item \emph{Say it:} \emph{neposlúshnyy}
  \item \emph{Say it:} \emph{neposlúshnyy}, once more
  \item \emph{Say it:} a sentence of your own with \emph{neposlúshnyy}
  \item \emph{Recall:} say the Russian for to obey, then the Russian for to be naughty, then say \emph{neposlúshnyy} again
\end{itemize}

\begin{tcolorbox}[breakable,colback=teal!4,colframe=teal!35!black,title={Before you move on}]
What does \ru{непослушный} mean? (\textquotedblleft{}Disobedient, naughty\textquotedblright{}.) Say it once more.
\end{tcolorbox}
